% ==================== 4.5 本章小结 ====================
\section{本章小结}
\label{sec:ch4_summary}

本章针对月面巡视器的动力学特性，系统研究了轮-壤交互力学理论和多体动力学建模方法，构建了巡视器动力学数字孪生模型。主要工作和结论如下：

（1）在轮-壤交互力学理论方面，回顾了经典地面力学理论，针对月面环境的特殊性，建立了考虑滑移特性的轮-壤相互作用模型。模型考虑了滑移沉陷、低重力效应等因素的影响，能够准确预测车轮的沉陷、牵引力和运动阻力。

（2）在多体动力学建模方面，分析了巡视器的拓扑结构，采用D-H参数法和Lagrange方法建立了巡视器的运动学和动力学方程。动力学模型考虑了轮-壤相互作用力的作用，能够描述巡视器在月面复杂地形下的运动行为。

（3）构建了巡视器动力学数字孪生仿真系统，实现了实时仿真。采用Runge-Kutta四阶方法进行数值积分，时间步长20 ms，满足实时性要求。

（4）通过地面模拟实验验证了模型的有效性。仿真结果与实验数据的误差在10\%以内，证明了模型的准确性。实时仿真结果表明，仿真时间比接近1，满足实时性要求。

本章建立的巡视器动力学数字孪生模型为后续路径规划、避障决策和运动控制提供了可靠的动力学支撑。
